% =====================================================================
\section{Representing Motion and Appearance}
\label{sec:setup}

\begin{figure}[tbp]
\centering
\includegraphics[width=\linewidth]{figures/splats/fig2_interleaved.png}
\caption{\textbf{Motion fingerprints (BFV splats) and RGB samples.} Each column is
one dataset. Within each group the top row shows an RGB sample and the row beneath
it the dataset's motion fingerprint. Each fingerprint pools the dataset's joint flow
vectors, position and displacement $[\hat x,\hat y,\Delta\hat x,\Delta\hat y]$ (the
bag-of-flow-vectors, BFV), into spatial clusters drawn as directional splats, where a
splat's hue gives its cluster's dominant flow direction (wheel) and its extent the
cluster spread. The top group holds the training sources and the bottom the target
benchmarks. The first three columns carry the datasets present in both sets
(FlyingThings, PointOdyssey, SPair) in a shared column, so the motion fingerprint is
visibly a property of the dataset that survives the split.}
\label{fig:splats}
\end{figure}

We study a wide range of correspondence and related tasks, across optical flow, dense geometric matching, and semantic keypoint matching. The datasets we consider are heterogeneous, some carrying dense per-pixel flow, others only sparse keypoint annotations. To bridge them, we map every dataset into a shared metric space and represent it as an empirical distribution $\mathcal{X}=\{\mathbf{x}_1,\dots,\mathbf{x}_n\}$ over per-sample feature vectors, where each $\mathbf{x}_i$ is a motion vector $\mathbf{f}_{\text{flow}}$ or an appearance vector $\mathbf{f}_{\text{app}}$ defined below.

\noindent\textbf{Motion representation.}
Because a motion field's spatial layout (radial expansion for ego-motion, vertical scatter for gravity dominated scenes) is signal that direction or magnitude alone discards~\cite{roth_opt_flow}, we encode motion as a joint distribution of position and displacement:
\begin{equation}
\mathbf{f}_{\text{flow}} = \bigl[\hat{x},\;\hat{y},\;\Delta\hat{x},\;\Delta\hat{y}\bigr],
\label{eq:bfv}
\end{equation}
We crop each image to a square ($H{=}W$), then normalize pixel coordinates to
$[-1,1]$ by image dimensions ($\hat{x} = 2x/W - 1$, $\hat{y} = 2y/H - 1$), and
displacements are expressed in the same units ($\Delta\hat{x} = 2\,\Delta x/W$,
$\Delta\hat{y} = 2\,\Delta y/H$; a displacement of $\pm1$ spans half the image).
For datasets where dense optical flow is available, the displacement
$(\Delta\hat{x},\Delta\hat{y})$ is the per-pixel flow vector; for datasets where only sparse
keypoint correspondences are annotated, the displacement is the keypoint
correspondence, treated as a sparse flow field. The encoding is resolution-invariant, free of spatial binning, and bounded by
construction (motion large enough to leave the frame has no valid correspondence and
drops out).  We call this the bag-of-flow-vectors (BFV) representation (Fig.~\ref{fig:splats}). A dataset's motion distribution $\mathcal{X}$ collects one such $\mathbf{f}_{\text{flow}}$ per pixel or keypoint, so its points are $\mathbf{x}_i = \mathbf{f}_{\text{flow}}$.

\noindent\textbf{Appearance representation.}
Motion is not the only plausible prior. Because two datasets can share flow statistics while differing in texture,
lighting, and object identity, appearance is the natural competing
hypothesis. Comparing deep features alone already gives competitive
zero-shot dense correspondence~\cite{sd_dino_tale_of_two_features}, so one could expect
a source that looks like the target to transfer well. We test
appearance and measure its predictive power against motion's on equal
footing. To do this we represent each dataset's appearance by its DINOv3 patch
features~\cite{simeoni2025dinov3}, which are learned self-supervised on natural images and naturally captures the texture, lighting, and object instance and position. We follow standard matching practice by whitening each dino patch embedding $\mathbf{z}_p$ (projecting
onto the leading principal components with $\mathbf{W}$) and normalizing to the unit sphere,
\begin{equation}
\mathbf{f}_{\text{app}} = \frac{\mathbf{W}(\mathbf{z}_p - \boldsymbol{\mu})}
{\lVert \mathbf{W}(\mathbf{z}_p - \boldsymbol{\mu}) \rVert_2}.
\label{eq:app}
\end{equation}
The appearance distribution $\mathcal{X}$ collects one such $\mathbf{f}_{\text{app}}$ per patch, so its points are $\mathbf{x}_i = \mathbf{f}_{\text{app}}$.

\noindent\textbf{Distribution Distances.}
We compare distributions in both representations with standard distributional distances (Chamfer,
FID~\cite{heusel2018ganstrainedtimescaleupdate}, and sliced Wasserstein-2~\cite{bonneel2015sliced}), and
further, decompose Chamfer distance into its two directed halves for a more nuanced comparison. Concretely, for distributions $A$ and $B$, the
directed distance $d_{A \to B}$ and the Chamfer
$d_{\mathrm{Cham}}$~\cite{fan2017pointcloud} that is the sum with the reverse
$d_{B \to A}$ (the same sum with the distributions swapped) are:
\begin{equation}
d_{A \to B} \;=\; \frac{1}{|A|}\sum_{a \in A}\; \min_{b \in B}\,
\lVert a - b \rVert_2^2,
\qquad
d_{\mathrm{Cham}} \;=\; d_{A \to B} + d_{B \to A}.
\label{eq:dirdist}
\end{equation}
Applied to a training-source distribution $S$ and a target-benchmark distribution $T$, the two
directions capture different failures. \textbf{Off-target mass} ($\dtb$) is large when
much of the training source's distribution sits far from anything the target benchmark
uses. \textbf{Coverage} ($\dbt$) is large when target-benchmark motions have no nearby
training-source example. Both are distances where
smaller is a closer match. We use
\emph{coverage} and \emph{off-target mass} for the two directions and Chamfer for
their sum throughout. The same distances apply unchanged to appearance
distributions.

Notably, these distances (Chamfer, FID, sliced Wasserstein-2) do not rely on a continuous density estimate. Therefore, we can compute them directly from the samples and adding more frames does not skew the estimate. KL and other density-based divergences, instead, rely on a $k$-NN density estimate that is sampling-sensitive and fails to converge on the heterogeneous datasets we study (see supplementary Sect.~\ref{S-supp:sampling}).
